% ==========================================================================
%  Arquivo de exemplo — formato LaTeX do QuestBank
%  Converter:  latex2questbank exemplo.tex -o exemplo.json
% ==========================================================================

\begin{questao}{000052}
  \meta{tipo}{objetiva}
  \meta{banca}{ENEM}
  \meta{ano}{2020}
  \meta{disciplina}{Física}
  \meta{topico}{Energia}
  \meta{conteudo}{Fontes de energia}
  \meta{assunto}{Energia eólica e impacto ambiental}
  \meta{dificuldade}{medio}
  \meta{tags}{vestibular, ENEM, energia}

  \enunciado{
    O uso de equipamentos elétricos custa dinheiro e libera carbono na
    atmosfera. Entretanto, diferentes usinas de energia apresentam custos
    econômicos e ambientais distintos.

    \imagem{grafico-energia.png}

    Em relação aos custos associados às fontes energéticas apresentadas,
    a energia obtida a partir do vento é
  }

  \begin{alternativas}
    \alt{A}{mais cara que a energia nuclear e emite maior quantidade de carbono.}
    \alt{B}{a segunda fonte mais cara e é livre de emissões de carbono.}
    \alt{C}{mais cara que a energia solar e ambas são livres de emissões.}
    \alt{D}{mais barata que as demais e emite grandes quantidades de carbono.}
    \alt{E}{a fonte que gera energia mais barata e livre de emissões.}
  \end{alternativas}

  \gabarito{E}
\end{questao}

% ----------------------- questão com fórmulas matemáticas -------------------

\begin{questao}{000073}
  \meta{tipo}{objetiva}
  \meta{banca}{FUVEST}
  \meta{ano}{2023}
  \meta{disciplina}{Física}
  \meta{topico}{Mecânica}
  \meta{conteudo}{Dinâmica}
  \meta{assunto}{Segunda Lei de Newton}
  \meta{dificuldade}{medio}

  \enunciado{
    Um corpo de massa $m = 2\,\text{kg}$ é submetido a uma força resultante
    $\vec{F} = 10\,\text{N}$. A aceleração do corpo, em m/s\textsuperscript{2}, é:

    $$a = \frac{F}{m}$$
  }

  \begin{alternativas}
    \alt{A}{$2\,\text{m/s}^2$}
    \alt{B}{$5\,\text{m/s}^2$}
    \alt{C}{$10\,\text{m/s}^2$}
    \alt{D}{$20\,\text{m/s}^2$}
  \end{alternativas}

  \gabarito{B}
\end{questao}

% ----------------------- questão discursiva ---------------------------------

\begin{questao}{000091}
  \meta{tipo}{discursiva}
  \meta{banca}{UNICAMP}
  \meta{ano}{2025}
  \meta{disciplina}{Física}
  \meta{topico}{Mecânica}
  \meta{conteudo}{Energia}
  \meta{assunto}{Energia Potencial Gravitacional}
  \meta{dificuldade}{dificil}

  \enunciado{
    Explique o conceito de \textbf{energia potencial gravitacional} e
    derive a expressão para a energia potencial gravitacional de um corpo
    de massa $m$ a uma altura $h$ da superfície terrestre.
  }

  \gabarito{$E_p = m g h$, onde $m$ é a massa, $g$ a aceleração gravitacional e $h$ a altura.}
\end{questao}

% ----------------------- questão adaptada (prefixo A-) ----------------------

\begin{questao}{A-000052}
  \meta{tipo}{objetiva}
  \meta{banca}{ENEM}
  \meta{ano}{2020}
  \meta{disciplina}{Física}
  \meta{topico}{Energia}
  \meta{conteudo}{Fontes de energia}
  \meta{assunto}{Energia eólica e impacto ambiental}
  \meta{dificuldade}{facil}
  \meta{tags}{adaptada, NEE, ENEM}

  \enunciado{
    O gráfico abaixo mostra o custo e a quantidade de carbono liberado por
    diferentes fontes de energia.

    \imagem{grafico-energia.png}

    \textbf{Analisando o gráfico, o que podemos afirmar sobre a energia
    eólica (vento)?}
  }

  \begin{alternativas}
    \alt{A}{É a energia mais cara do gráfico e não emite carbono.}
    \alt{B}{É a energia mais barata do gráfico e não emite carbono.}
    \alt{C}{É mais barata que a nuclear, mas emite muito carbono.}
  \end{alternativas}

  \gabarito{B}
\end{questao}
